\section{Git and GitHub Foundations}


\begin{frame}{From Bash to Version Control}
\framesubtitle{The terminal operates files; Git manages their engineering history}
\footnotesize
\begin{center}
\colorbox{um6p-color-dark-green!18}{\parbox[c][0.92cm][c]{2.15cm}{\centering\textbf{Ubuntu}\\\scriptsize Completed}}
\quad$\Longrightarrow$\quad
\colorbox{um6p-color-dark-green!18}{\parbox[c][0.92cm][c]{2.15cm}{\centering\textbf{Terminal and Bash}\\\scriptsize Completed}}
\quad$\Longrightarrow$\quad
\colorbox{um6p-color-light-blue}{\parbox[c][1.02cm][c]{2.25cm}{\centering\textbf{Git and GitHub}\\\scriptsize Current foundation}}
\end{center}
\begin{block}{New Engineering Question}
How can several students modify the same robotics project while preserving its history, identifying each contribution, and keeping the GitHub repository reproducible?
\end{block}
\begin{alertblock}{Central Principle}
Git is not a replacement for careful engineering. It makes careful work visible, traceable, reviewable, and recoverable.
\end{alertblock}
\end{frame}

\begin{frame}{Why Robotics Needs Version Control}
\framesubtitle{Robotics systems evolve through many connected files and contributors}
\footnotesize
\begin{columns}[T,onlytextwidth]
\begin{column}{0.48\textwidth}
\begin{block}{Without Version Control}
\begin{itemize}\setlength{\itemsep}{0.2em}
\item files become \texttt{final}, \texttt{final2}, and \texttt{final\_fixed};
\item working configurations are overwritten;
\item changes cannot be attributed or reviewed;
\item two contributors may replace each other's work.
\end{itemize}
\end{block}
\end{column}
\begin{column}{0.48\textwidth}
\begin{exampleblock}{With Version Control}
\begin{itemize}\setlength{\itemsep}{0.2em}
\item each meaningful change becomes a recorded checkpoint;
\item experiments can be isolated on branches;
\item collaborators synchronize controlled changes;
\item earlier states remain available for investigation.
\end{itemize}
\end{exampleblock}
\end{column}
\end{columns}
\begin{alertblock}{Project Connection}
A change to a ROS 2 node may depend on a launch file, model, environment variable, or README instruction. Version control preserves these relationships as the system evolves.
\end{alertblock}
\end{frame}

\begin{frame}{What Is Version Control?}
\framesubtitle{A structured history of a project's meaningful states}
\footnotesize
\begin{center}
Working files $\longrightarrow$ Selected changes $\longrightarrow$ Recorded checkpoint $\longrightarrow$ Shared history
\end{center}
\begin{block}{Intuitive Definition}
Version control records how a collection of files changes over time. It helps engineers answer: \emph{What changed? Who changed it? Why was it changed? When was it changed?}
\end{block}

\end{frame}






\begin{frame}{Git and GitHub}
\framesubtitle{A local version-control tool and an online collaboration platform}
\footnotesize
\begin{columns}[T,onlytextwidth]
\begin{column}{0.47\textwidth}
\begin{block}{Git}
\begin{itemize}\setlength{\itemsep}{0.18em}
\item distributed version-control software;
\item runs on the local computer;
\item records commits, branches, and merges;
\item can work locally without GitHub.
\end{itemize}
\end{block}
\end{column}
\begin{column}{0.47\textwidth}
\begin{exampleblock}{GitHub}
\begin{itemize}\setlength{\itemsep}{0.18em}
\item online hosting and collaboration platform;
\item stores shared Git repositories;
\item provides browsing, issues, reviews, and documentation;
\item connects contributors through remote repositories.
\end{itemize}
\end{exampleblock}
\end{column}
\end{columns}
\end{frame}



\begin{frame}{Git and GitHub}

\footnotesize

\begin{columns}[T,onlytextwidth]

    % --- LEFT COLUMN: Git (The Camera) ---
    \begin{column}{0.48\textwidth}
        \begin{block}{Git: The Local Camera}
            Think of Git as a digital camera for your project.
            \begin{itemize}
                \setlength{\itemsep}{0.2em}
                \item When you write a good piece of code, you take a "snapshot" (\textbf{commit}).
                \item If your next experiment breaks the drone, you can simply rewind to an older snapshot.
                \item \textbf{Limitation:} These snapshots live \textit{only} on your physical laptop. If the laptop breaks, the code is gone.
            \end{itemize}
        \end{block}
    \end{column}

    % --- RIGHT COLUMN: GitHub (The Cloud) ---
    \begin{column}{0.48\textwidth}
        \begin{exampleblock}{GitHub: The Shared Gallery}
            Think of GitHub as an online photo album.
            \begin{itemize}
                \setlength{\itemsep}{0.2em}
                \item You upload your local snapshots to the internet (\textbf{push}).
                \item Your 9 teammates can browse your code and download it to their own laptops (\textbf{pull}).
                \item \textbf{Advantage:} It acts as the ultimate backup. It is the "single source of truth" for the team.
            \end{itemize}
        \end{exampleblock}
    \end{column}

\end{columns}

\pause
\vspace{0.25cm}

% --- BOTTOM: The Daily Workflow ---
\begin{alertblock}{The Daily Project Workflow}
    \centering
    \normalsize
    \textbf{1. Write Code} 
    \quad $\longrightarrow$ \quad 
    \textbf{2. Snap Photo} (\texttt{git commit}) 
    \quad $\longrightarrow$ \quad 
    \textbf{3. Upload} (\texttt{git push})
\end{alertblock}

\end{frame}

\begin{frame}{Essential Git Vocabulary}
\framesubtitle{A project-centered mental model before using commands}
\scriptsize
\renewcommand{\arraystretch}{1.16}
\begin{tabularx}{\textwidth}{>{\bfseries\color{um6p-color-dark-blue}}p{0.18\textwidth} X}
\toprule
Repository & Project files together with the Git history and metadata. \\
Commit & A recorded snapshot of selected changes with author, time, and message. \\
Branch & A movable line of development used to isolate related work. \\
Remote & A named connection to another repository, commonly \texttt{origin} on GitHub. \\
Clone & A local copy of a repository, including its history and remote connection. \\
Pull & Retrieve remote changes and integrate them into the current local branch. \\
Push & Publish local commits to a remote branch. \\
Merge & Combine histories from different lines of development. \\
\texttt{.gitignore} & Rules describing untracked files that Git should normally ignore. \\
\bottomrule
\end{tabularx}
\end{frame}



\begin{frame}[fragile]{Before Cloning: Verify the Destination}
\framesubtitle{Clone the repository once and into the intended parent directory}
\footnotesize
\begin{block}{Pre-Clone Checks}
\begin{verbatim}
cd ~
pwd
ls -la
which git
git --version
\end{verbatim}
\end{block}
\begin{exampleblock}{Question Before Execution}
Will \texttt{git clone} create \texttt{quard\_um6p\_intership} inside the current directory, or does a folder with that name already exist?
\end{exampleblock}
\begin{alertblock}{Common Mistake}
Do not clone the repository repeatedly inside itself or inside another accidental clone. If the target folder already exists, inspect it before taking action.
\end{alertblock}
\end{frame}

\begin{frame}{Access the Internship Repository on GitHub}
\framesubtitle{Understand the platform before copying commands}
\footnotesize
\begin{block}{Repository Address}
\centering
\url{https://github.com/abdtnourji/quard_um6p_intership}
\end{block}
\begin{columns}[T,onlytextwidth]
\begin{column}{0.48\textwidth}
\begin{exampleblock}{Inspect on GitHub}
\begin{itemize}\setlength{\itemsep}{0.18em}
\item repository name and owner;
\item default branch;
\item README and documentation;
\item folders and recent commits;
\item access permissions.
\end{itemize}
\end{exampleblock}
\end{column}
\begin{column}{0.48\textwidth}
\begin{block}{Obtain the URL}
Select \textbf{Code}, choose the documented authentication method, and copy the repository URL exactly. HTTPS is the simplest starting point when no SSH key workflow has been configured.
\end{block}
\end{column}
\end{columns}
\end{frame}

\begin{frame}[fragile]{Clone the Internship Repository}
\framesubtitle{Create one local working copy with history and remote configuration}
\footnotesize
\begin{block}{Purpose and Syntax}
\begin{verbatim}
git clone <repository-url>
\end{verbatim}
\end{block}
\begin{exampleblock}{Internship Command}
\begin{verbatim}
cd ~
git clone https://github.com/abdtnourji/quard_um6p_intership.git
cd quard_um6p_intership
\end{verbatim}
\end{exampleblock}
\begin{block}{What the Command Creates}
The project directory, working files, commit history, branches known at clone time, and a remote normally named \texttt{origin}.
\end{block}
% \begin{alertblock}{If Authentication Is Requested}
% Do not enter a GitHub account password as though it were a reusable Git password. Follow the repository's approved HTTPS credential or SSH-key procedure.
% \end{alertblock}
\end{frame}

\begin{frame}[fragile]{Verify That the Clone Is Correct}
\framesubtitle{A successful-looking download still requires engineering checks}
\footnotesize
\begin{exampleblock}{Verification Sequence}
\begin{verbatim}
cd ~/quard_um6p_intership
pwd
ls -la
git status
git branch
git remote -v
git log --oneline -5
\end{verbatim}
\end{exampleblock}
\begin{block}{Expected Evidence}
The path identifies the intended repository, \texttt{.git} exists, Git recognizes a branch, the working tree is initially clean, \texttt{origin} points to the expected GitHub URL, and recent commits are visible.
\end{block}
% \begin{alertblock}{Repository Test}
% If \texttt{git status} reports ``not a git repository,'' verify the current directory and move to the cloned project root.
% \end{alertblock}
\end{frame}

\begin{frame}[fragile]{Explore Before Modifying}
\framesubtitle{Understand the project structure and follow its documentation}
\footnotesize
\begin{columns}[T,onlytextwidth]
\begin{column}{0.48\textwidth}
\begin{block}{Filesystem View}
\begin{verbatim}
ls -la
find . -maxdepth 2 -type d | sort
find . -name "README*"
find . -name "package.xml"
\end{verbatim}
\end{block}
\end{column}
\begin{column}{0.48\textwidth}
\begin{exampleblock}{Git-Aware View}
\begin{verbatim}
git status
git branch
git log --oneline --graph -8
git remote -v
\end{verbatim}
\end{exampleblock}
\end{column}
\end{columns}
\begin{alertblock}{Professional Rule}
The README is part of the engineering interface. Read setup instructions, folder responsibilities, prerequisites, and known limitations before executing the pipeline.
\end{alertblock}
\end{frame}

\begin{frame}[fragile]{Use \texttt{git status} as the Dashboard}
\framesubtitle{Ask Git what it sees before and after every important action}
\footnotesize
\begin{block}{Purpose and Syntax}
\begin{verbatim}
git status
git status --short
\end{verbatim}
\end{block}
\begin{exampleblock}{Interpret the Categories}
\begin{itemize}\setlength{\itemsep}{0.18em}
\item \textbf{untracked}: Git has not been asked to include the file;
\item \textbf{modified}: tracked content differs from the last commit;
\item \textbf{staged}: content is selected for the next commit;
\item \textbf{clean}: no uncommitted change is detected.
\end{itemize}
\end{exampleblock}
\begin{alertblock}{Habit}
Run \texttt{git status} before editing, before staging, before committing, before pulling, and before pushing. It provides context, not just success or failure.
\end{alertblock}
\end{frame}

\begin{frame}[fragile]{Create a Small, Safe Change}
% \framesubtitle{Practise the workflow with student-owned documentation}
\footnotesize
\begin{exampleblock}{Project Exercise}
\begin{verbatim}
cd ~/quard_um6p_intership
mkdir -p student_work
echo "Git workflow observations" > student_work/git_notes.txt
git status
\end{verbatim}
\end{exampleblock}
\begin{block}{Engineering Intention}
This creates a visible, low-risk change without modifying flight logic, simulation models, or shared configuration.
\end{block}
\begin{alertblock}{Repository Policy}
If the project defines a specific student folder, naming convention, or contribution process, follow that documentation instead of inventing a parallel structure.
\end{alertblock}
\end{frame}

\begin{frame}[fragile]{Inspect Changes with \texttt{git diff}}
\framesubtitle{Review content before selecting it for a commit}
\footnotesize
\begin{block}{Purpose and Syntax}
\begin{verbatim}
git diff              # unstaged tracked changes
git diff --staged     # changes selected for commit
\end{verbatim}
\end{block}
\begin{exampleblock}{Project Workflow}
\begin{verbatim}
echo "Repository inspected from the project root." \
  >> student_work/git_notes.txt
git diff
git status
\end{verbatim}
\end{exampleblock}
\begin{alertblock}{Important Detail}
An untracked file may not appear in ordinary \texttt{git diff}. Use \texttt{git status} to discover it, then stage it before reviewing with \texttt{git diff --staged}.
\end{alertblock}
\end{frame}

\begin{frame}[fragile]{Stage Only the Intended Change}
\framesubtitle{Build the next commit deliberately}
\footnotesize
\begin{block}{Purpose and Syntax}
\begin{verbatim}
git add <file-or-directory>
\end{verbatim}
\end{block}
\begin{exampleblock}{Preferred Beginner Workflow}
\begin{verbatim}
git add student_work/git_notes.txt
git status
git diff --staged
\end{verbatim}
\end{exampleblock}
\begin{alertblock}{Why Not Always Use \texttt{git add .}?}
Staging everything may include logs, generated files, credentials, unrelated edits, or large assets. Add explicit paths until you can explain every staged change.
\end{alertblock}
\end{frame}

\begin{frame}[fragile]{Create a Meaningful Commit}
\framesubtitle{Record one coherent engineering change with a useful message}
\footnotesize
\begin{block}{Purpose and Syntax}
\begin{verbatim}
git commit -m "concise description of the change"
\end{verbatim}
\end{block}
\begin{columns}[T,onlytextwidth]
\begin{column}{0.48\textwidth}
\begin{exampleblock}{Useful Message}
\begin{verbatim}
git commit -m \
  "docs: add student Git workflow notes"
\end{verbatim}
It states the area and result.
\end{exampleblock}
\end{column}
\begin{column}{0.48\textwidth}
\begin{alertblock}{Weak Messages}
\texttt{update}, \texttt{changes}, \texttt{final}, and \texttt{fix stuff} do not explain engineering intent.
\end{alertblock}
\end{column}
\end{columns}
\begin{block}{After Committing}
Run \texttt{git status} and \texttt{git log --oneline -5}. The commit exists locally until it is pushed.
\end{block}
\end{frame}

\begin{frame}[fragile]{Read the Project History}
\framesubtitle{Commits form an engineering narrative, not merely a backup list}
\footnotesize
\begin{block}{Purpose and Useful Views}
\begin{verbatim}
git log
git log --oneline
git log --oneline --graph --decorate --all
\end{verbatim}
\end{block}
\begin{exampleblock}{Questions the History Can Answer}
\begin{itemize}\setlength{\itemsep}{0.18em}
\item Which change introduced a launch instruction?
\item When was a model or configuration updated?
\item Which branch contains a student's contribution?
\item What commit preceded a newly observed problem?
\end{itemize}
\end{exampleblock}
% \begin{alertblock}{Quality Signal}
% Small coherent commits with descriptive messages are easier to review, merge, test, and recover than one large unexplained commit.
% \end{alertblock}
\end{frame}

\begin{frame}[fragile]{Branches Isolate Lines of Work}
\framesubtitle{Experiment without immediately changing the shared main line}
\footnotesize
\begin{block}{Inspect and Create}
\begin{verbatim}
git branch
git switch -c docs/student-git-notes
\end{verbatim}
\texttt{git branch} lists branches. \texttt{git switch -c} creates and enters a new branch.
\end{block}
\begin{exampleblock}{Move Between Existing Branches}
\begin{verbatim}
git switch main
git switch docs/student-git-notes
\end{verbatim}
\end{exampleblock}
\begin{alertblock}{Before Switching}
Check \texttt{git status}. Uncommitted changes may follow you, conflict with the destination, or prevent the switch. Branch names should describe the task, not the contributor's mood.
\end{alertblock}
\end{frame}

\begin{frame}{What Does Merge Mean?}
\framesubtitle{Combine compatible histories after review and verification}
\footnotesize
\begin{center}
\texttt{main}: A $\longrightarrow$ B $\longrightarrow$ C

\vspace{0.15cm}
\hspace{1.0cm}$\searrow$ D $\longrightarrow$ E \quad \texttt{feature branch}

\vspace{0.15cm}
$\Downarrow$

\texttt{main}: A $\longrightarrow$ B $\longrightarrow$ C $\longrightarrow$ M
\end{center}
\begin{block}{Intuition}
A merge combines changes from different lines of development. If both lines modify incompatible parts of the same content, Git asks an engineer to resolve the conflict.
\end{block}
\begin{alertblock}{Responsibility}
A conflict is not Git failing. It is Git refusing to guess which engineering intention is correct. Resolve it by understanding both changes, then test the combined result.
\end{alertblock}
\end{frame}

\begin{frame}[fragile]{Inspect the Remote Connection}
\framesubtitle{Know where pulls come from and where pushes go}
\footnotesize
\begin{block}{Purpose and Syntax}
\begin{verbatim}
git remote -v
\end{verbatim}
\end{block}
\begin{exampleblock}{Expected Interpretation}
\texttt{origin} is the conventional name created by cloning. The output normally shows fetch and push URLs for the GitHub repository.
\end{exampleblock}
\begin{alertblock}{Wrong-Remote Prevention}
Before your first push, confirm that the owner and repository name are correct and that you have authorized contribution access. Do not replace \texttt{origin} randomly to bypass a permission problem.
\end{alertblock}
\end{frame}

\begin{frame}[fragile]{Synchronize Before Publishing}
\framesubtitle{Bring in remote work before sending your own commits}
\footnotesize
\begin{block}{Safe Beginner Sequence}
\begin{verbatim}
git status
git branch
git pull origin main
\end{verbatim}
\end{block}
\begin{exampleblock}{Meaning}
\texttt{git pull} retrieves remote changes and integrates them into the current local branch. Read the output to determine whether the branch was already current, fast-forwarded, or requires conflict resolution.
\end{exampleblock}
\begin{alertblock}{Before Pulling}
Start from a clean, understood working state. If local edits are unfinished, commit them appropriately or follow the instructor's recovery procedure rather than forcing the pull.
\end{alertblock}
\end{frame}

\begin{frame}[fragile]{Publish Local Commits with \texttt{git push}}
\framesubtitle{Synchronize verified history to GitHub}
\footnotesize
\begin{block}{Purpose and Syntax}
\begin{verbatim}
git push <remote> <branch>
\end{verbatim}
\end{block}
\begin{exampleblock}{Examples}
\begin{verbatim}
git push origin main
# For a new contribution branch:
git push -u origin docs/student-git-notes
\end{verbatim}
\end{exampleblock}
\begin{alertblock}{Rejected Push}
A non-fast-forward rejection usually means the remote contains commits missing locally. Do not force-push. Retrieve and understand upstream changes, resolve any conflict, test, and then push normally.
\end{alertblock}
\end{frame}

\begin{frame}{The Everyday Internship Workflow}
\framesubtitle{Observe, synchronize, modify, review, record, and share}
\scriptsize
\begin{enumerate}\setlength{\itemsep}{0.28em}
\item \textbf{Enter and inspect:} move to the repository, then run \texttt{git status} and \texttt{git branch}.
\item \textbf{Synchronize:} pull the approved branch before beginning new work.
\item \textbf{Modify:} make one focused code, configuration, test, or documentation change.
\item \textbf{Review:} inspect \texttt{git status} and \texttt{git diff}.
\item \textbf{Stage:} add only the files intended for the next checkpoint.
\item \textbf{Verify staged content:} inspect \texttt{git diff --staged}.
\item \textbf{Commit:} use a concise message that explains the result.
\item \textbf{Synchronize again:} pull if necessary, resolve carefully, test, and push.
\item \textbf{Confirm on GitHub:} verify the branch, commit, and documentation online.
\end{enumerate}
\begin{alertblock}{Team Rule}
The exact branch and review policy is defined by the internship repository. Follow the documented collaboration workflow rather than assuming every contributor pushes directly to \texttt{main}.
\end{alertblock}
\end{frame}


\begin{frame}{Protecting the Project: Forking}
\framesubtitle{Creating your own workspace without impacting the main repository}

\footnotesize

\begin{columns}[T,onlytextwidth]

    % --- LEFT COLUMN: The Concept ---
    \begin{column}{0.48\textwidth}
        \begin{alertblock}{The "Golden Copy" Rule}
            The supervisor's main GitHub repository is the \textbf{Golden Copy}. If 10 students push their experimental code directly into it at the same time, the project will instantly break.
        \end{alertblock}

        \begin{block}{The Solution: "Forking"}
            A \textbf{Fork} is a completely separate, personal copy of the repository that lives on \textit{your} GitHub account. 
            \begin{itemize}
                \setlength{\itemsep}{0.2em}
                \item You can freely modify it.
                \item You can completely break it.
                \item It \textbf{never} affects the supervisor's original code.
            \end{itemize}
        \end{block}
    \end{column}

    % --- RIGHT COLUMN: The Workflow ---
    \begin{column}{0.48\textwidth}
        \begin{exampleblock}{Your Safe Setup Workflow}
            Instead of cloning the main project directly, you will:
            \begin{enumerate}
                \setlength{\itemsep}{0.3em}
                \item \textbf{Fork (In Browser):} Go to the supervisor's GitHub page and click the "Fork" button to create your copy.
                \item \textbf{Clone (In Terminal):} Run \texttt{git clone} using the URL of \textit{your} new fork, not the supervisor's URL.
                \item \textbf{Push (In Terminal):} When you type \texttt{git push}, your code goes safely to your personal cloud repository.
            \end{enumerate}
        \end{exampleblock}
    \end{column}

\end{columns}



\end{frame}



\begin{frame}[fragile]{Use \texttt{.gitignore} Deliberately}
\framesubtitle{Keep generated, local, and sensitive artifacts out of history}

\footnotesize

\begin{columns}[T,onlytextwidth]

    % --- LEFT COLUMN: Theory & Warnings ---
    \begin{column}{0.52\textwidth}
        \begin{block}{The Purpose}
            A \texttt{.gitignore} file contains patterns telling Git to ignore specific untracked files, preventing them from accidentally being committed and pushed to GitHub.
        \end{block}
        
        \begin{alertblock}{Important Boundaries}
            \begin{itemize}
                \setlength{\itemsep}{0.3em}
                \item \textbf{Too Late:} Ignoring a file does \textit{not} remove it from history if it was already committed previously.
                \item \textbf{Security:} Never rely on \texttt{.gitignore} as your only protection for credentials. Never commit passwords.
                \item \textbf{Bloat:} Use it to prevent uploading massive ROS 2 logs or compiled build files that crash repositories.
            \end{itemize}
        \end{alertblock}
    \end{column}

    % --- RIGHT COLUMN: The Code Example ---
    \begin{column}{0.44\textwidth}
        \begin{exampleblock}{Typical Robotics Patterns}
\begin{verbatim}
# Python cache
__pycache__/
*.pyc
# Local ROS 2 workspaces
build/
install/
log/
# Local environments
.venv/
# Sensitive secrets
.env
*.pem
\end{verbatim}
        \end{exampleblock}
    \end{column}

\end{columns}

\end{frame}

\begin{frame}{Never Commit Secrets}
\framesubtitle{Repository history can outlive the working file}
\footnotesize
\begin{alertblock}{Do Not Commit}
\begin{itemize}\setlength{\itemsep}{0.20em}
\item passwords, access tokens, API keys, or private SSH keys;
\item private certificates or machine-specific credentials;
\item confidential datasets or personal information;
\item environment files containing secret values.
\end{itemize}
\end{alertblock}
\begin{block}{If a Secret Is Exposed}
Stop sharing it, notify the responsible instructor or repository maintainer, revoke or rotate the credential, and follow the approved history-cleaning procedure. Merely deleting the file in a later commit may leave the secret in earlier history.
\end{block}
\begin{exampleblock}{Better Pattern}
Commit a documented template such as \texttt{.env.example} with placeholder names, then keep real secret values outside version control.
\end{exampleblock}
\end{frame}

\begin{frame}{Keep the Repository Engineering-Friendly}
\framesubtitle{Version the source of truth, not every generated artifact}
\footnotesize
\begin{columns}[T,onlytextwidth]
\begin{column}{0.48\textwidth}
\begin{block}{Good Candidates}
\begin{itemize}\setlength{\itemsep}{0.18em}
\item source code and launch files;
\item human-maintained configuration;
\item small test fixtures;
\item README files and reports;
\item scripts needed to reproduce outputs.
\end{itemize}
\end{block}
\end{column}
\begin{column}{0.48\textwidth}
\begin{exampleblock}{Review Before Adding}
\begin{itemize}\setlength{\itemsep}{0.18em}
\item build and cache directories;
\item large videos, models, bags, or datasets;
\item temporary logs and screenshots;
\item machine-specific settings;
\item files obtainable from documented sources.
\end{itemize}
\end{exampleblock}
\end{column}
\end{columns}
\begin{alertblock}{Reason}
Large binary files make cloning, reviewing, and history management harder. Use the repository's approved storage and release strategy for necessary large assets.
\end{alertblock}
\end{frame}

\begin{frame}[fragile]{Recover: Staged the Wrong File}
\framesubtitle{Correct the selection without destroying the working change}
\footnotesize
\begin{exampleblock}{Situation}
You staged a file that should not be part of the next commit.
\begin{verbatim}
git status
git restore --staged path/to/file
git status
\end{verbatim}
\end{exampleblock}
\begin{block}{Effect}
The file is removed from the staging area, but its working-tree modification is preserved for later review.
\end{block}
\begin{alertblock}{Do Not Guess Recovery Commands}
Always inspect \texttt{git status} before and after recovery. Commands involving \texttt{--hard}, forced pushes, or history rewriting are outside the beginner workflow unless supervised.
\end{alertblock}
\end{frame}

\begin{frame}[fragile]{Recover: Discard an Unwanted Local Edit}
\framesubtitle{Use only when the last committed version is truly the desired state}
\footnotesize
\begin{exampleblock}{For One Tracked, Unstaged File}
\begin{verbatim}
git diff -- path/to/file
git restore path/to/file
git status
\end{verbatim}
\end{exampleblock}
\begin{alertblock}{Irreversible Working-Tree Loss}
\texttt{git restore path/to/file} discards the selected uncommitted edit. Copy valuable experimental content elsewhere or commit it on an appropriate branch before restoring.
\end{alertblock}
\begin{block}{Safer Question}
Do I want to remove the file from the next commit, or destroy the actual edit? These are different intentions and require different commands.
\end{block}
\end{frame}

\begin{frame}[fragile]{Recover: Push Rejected or Pull Conflicted}
\framesubtitle{Preserve evidence and resolve the real synchronization problem}
\scriptsize
\begin{columns}[T,onlytextwidth]
\begin{column}{0.48\textwidth}
\begin{block}{Push Rejected}
\begin{verbatim}
git status
git branch
git remote -v
git pull origin main
\end{verbatim}
Read the output, integrate approved remote work, test, and push again. Never use force as the automatic response.
\end{block}
\end{column}
\begin{column}{0.48\textwidth}
\begin{exampleblock}{Merge Conflict}
\begin{enumerate}\setlength{\itemsep}{0.15em}
\item run \texttt{git status};
\item open each conflicted file;
\item understand both alternatives;
\item remove conflict markers after deciding;
\item test the integrated result;
\item stage and commit the resolution.
\end{enumerate}
\end{exampleblock}
\end{column}
\end{columns}
\begin{alertblock}{Escalation Rule}
If the correct engineering intention is unclear, stop and ask the file owner or instructor. A syntactically resolved conflict can still create incorrect robot behavior.
\end{alertblock}
\end{frame}

\begin{frame}[fragile]{Two Computers, One Shared Repository}
% \framesubtitle{Synchronize through commits, not by copying folders manually}
\footnotesize
\begin{block}{At the Computer Where Work Was Completed}
\begin{verbatim}
git status
git add <intended-files>
git commit -m "describe the completed change"
git pull origin main
git push origin main
\end{verbatim}
\end{block}
\begin{exampleblock}{At the Other Computer}
\begin{verbatim}
cd ~/quard_um6p_intership
git status
git pull origin main
\end{verbatim}
\end{exampleblock}
\begin{alertblock}{Precondition}
Each computer must be on the intended repository and branch. Uncommitted local edits must be reviewed before pulling, or they may conflict with the incoming work.
\end{alertblock}
\end{frame}

\begin{frame}[fragile]{End-to-End Practice}
\framesubtitle{Obtain, modify, record, synchronize, and verify}
\scriptsize
\begin{exampleblock}{Guided Workflow}
\begin{verbatim}
cd ~
git clone https://github.com/abdtnourji/quard_um6p_intership.git
cd quard_um6p_intership

git status
git remote -v
git log --oneline -5

mkdir -p student_work
echo "Git practice completed" > student_work/git_practice.txt

git status
git add student_work/git_practice.txt
git diff --staged
git commit -m "docs: add Git practice evidence"
git pull origin main
git push origin main
\end{verbatim}
\end{exampleblock}

\end{frame}


\begin{frame}{References}
\framesubtitle{Primary documentation for continued practice}
\scriptsize
\begin{thebibliography}{9}
\bibitem{gitdocs}
Git Project, \emph{Git Documentation}.\\
\url{https://git-scm.com/doc}

\bibitem{gittutorial}
Git Project, \emph{A Tutorial Introduction to Git}.\\
\url{https://git-scm.com/docs/gittutorial}

\bibitem{githubbasics}
GitHub Docs, \emph{Git Basics}.\\
\url{https://docs.github.com/en/get-started/git-basics}

\bibitem{githubclone}
GitHub Docs, \emph{Cloning a Repository}.\\
\url{https://docs.github.com/en/repositories/creating-and-managing-repositories/cloning-a-repository}

\bibitem{githubpush}
GitHub Docs, \emph{Pushing Commits to a Remote Repository}.\\
\url{https://docs.github.com/en/get-started/using-git/pushing-commits-to-a-remote-repository}
\end{thebibliography}
\end{frame}




\begin{frame}{Transition to the Technical Foundations}
\framesubtitle{One foundation completed, the next one begins}

% \footnotesize

% ------------------------------------------------------------------
% CURRENT PROGRESS
% ------------------------------------------------------------------
\begin{center}
    \textbf{\color{um6p-dark}Internship Learning Roadmap}
\end{center}

\vspace{0.15cm}

\begin{center}
\renewcommand{\arraystretch}{1.15}
\begin{tabular}{c@{\hspace{0.10cm}}c@{\hspace{0.10cm}}c@{\hspace{0.10cm}}c@{\hspace{0.10cm}}c}

% Completed foundation
\colorbox{um6p-color-dark-green!20}{
    \parbox[c][0.85cm][c]{1.65cm}{
        \centering
        \scriptsize\bfseries
        Ubuntu\\
        \textcolor{um6p-color-dark-green}{Completed}
    }
}


& $\longrightarrow$ &

% Completed foundation
\colorbox{um6p-color-dark-green!20}{
    \parbox[c][1.05cm][c]{2.15cm}{
        \centering
        \scriptsize\bfseries
        Terminal\\
        and Bash\\[-0.05cm]
        \textcolor{um6p-color-dark-green}{Completed}
    }
}

& $\longrightarrow$ &

% Upcoming foundations
\colorbox{um6p-color-dark-green!20}{
    \parbox[c][0.85cm][c]{1.65cm}{
        \centering
        \scriptsize
        Git and\\
        GitHub \\[-0.05cm]
        \textcolor{um6p-color-dark-green}{Completed}
    }
}

\\[0.25cm]

\colorbox{um6p-color-light-blue}{
    \parbox[c][0.85cm][c]{1.65cm}{
        \centering
        \scriptsize
        Python for\\
        Robotics \\ [-0.05cm]
        \textcolor{um6p-dark}{Current}

    }
}

& $\longrightarrow$ &

\colorbox{black!7}{
    \parbox[c][0.85cm][c]{2.15cm}{
        \centering
        \scriptsize
        ROS 2
    }
}

& $\longrightarrow$ &

\colorbox{black!7}{
    \parbox[c][0.85cm][c]{1.65cm}{
        \centering
        \scriptsize
        Gazebo\\
        and PX4
    }
}

\end{tabular}

\vspace{0.15cm}

\colorbox{black!7}{
    \parbox[c][0.75cm][c]{2.1cm}{
        \centering
        \scriptsize
        AI Integration
    }
}
\end{center}


\end{frame}



% -----------------------------------------------------------------------------
\begin{frame}{Live Git Demonstration}
\framesubtitle{From the GitHub repository to a documented contribution}

\footnotesize

\begin{block}{Objective}
During this demonstration, we will reproduce the basic workflow used
throughout the internship:
\end{block}

\begin{center}
\scriptsize
\textbf{Clone}
$\longrightarrow$
\textbf{Inspect}
$\longrightarrow$
\textbf{Create a branch}
$\longrightarrow$
\textbf{Modify}
$\longrightarrow$
\textbf{Review}
$\longrightarrow$
\textbf{Commit}
$\longrightarrow$
\textbf{Push}
\end{center}

\pause

\begin{exampleblock}{Safe Practice Change}
We will create one small documentation file inside a dedicated student
folder. We will not modify the drone, PX4, Gazebo, or ROS 2 source code.
\end{exampleblock}

\begin{alertblock}{Important}
Execute one command at a time. Read its output before continuing to the
next step.
\end{alertblock}

\end{frame}

% -----------------------------------------------------------------------------
\begin{frame}{Before Starting}
\framesubtitle{Fork the repository if direct contribution is not permitted}

\footnotesize

\begin{block}{Recommended Student Workflow}
Open the internship repository in a web browser:

\begin{center}
\url{https://github.com/abdtnourji/quard_um6p_intership}
\end{center}

Then select \textbf{Fork} to create a personal copy under your own
GitHub account.
\end{block}

\pause

\begin{exampleblock}{Why Use a Fork?}
A fork gives each student a personal GitHub copy of the project.

Students can create branches, push experiments, and practise the Git
workflow without directly changing the supervisor's main repository.
\end{exampleblock}

\begin{alertblock}{Before Cloning}
Copy the URL of your own fork. Replace
\texttt{YOUR-GITHUB-USERNAME} in the following slides with your actual
GitHub username.
\end{alertblock}

\end{frame}

% -----------------------------------------------------------------------------
\begin{frame}[fragile]{Step 1: Check That Git Is Available}
\framesubtitle{Verify the tool before using it}

\footnotesize

\begin{block}{Execute}
\begin{verbatim}
which git
git --version
\end{verbatim}
\end{block}

\begin{exampleblock}{Expected Observation}
The first command should display the location of the Git executable.

The second command should display an installed Git version.
\end{exampleblock}

\begin{alertblock}{If the Command Is Not Found}
Do not continue by guessing. Verify that Git is installed using the
installation procedure documented for Ubuntu 22.04.
\end{alertblock}

\end{frame}

% -----------------------------------------------------------------------------
\begin{frame}[fragile]{Step 2: Prepare the Parent Directory}
\framesubtitle{Know where the repository will be created}

\footnotesize

\begin{block}{Execute}
\begin{verbatim}
cd ~
pwd
ls -la
\end{verbatim}
\end{block}

\begin{exampleblock}{Interpretation}
\begin{itemize}
    \setlength{\itemsep}{0.20em}
    \item \texttt{cd \textasciitilde} moves to the home directory;
    \item \texttt{pwd} confirms the current location;
    \item \texttt{ls -la} shows existing files and directories.
\end{itemize}
\end{exampleblock}

\begin{alertblock}{Check Before Cloning}
If a directory named \texttt{quard\_um6p\_intership} already exists,
do not clone another copy over it. Inspect the existing directory first.
\end{alertblock}

\end{frame}

% -----------------------------------------------------------------------------
\begin{frame}[fragile]{Step 3: Clone Your Fork}
\framesubtitle{Create a local working copy of the GitHub repository}

\footnotesize

\begin{block}{Execute}
\begin{verbatim}
git clone https://github.com/YOUR-GITHUB-USERNAME/\
quard_um6p_intership.git
\end{verbatim}
\end{block}

\begin{exampleblock}{What \texttt{git clone} Creates}
The command downloads:

\begin{itemize}
    \setlength{\itemsep}{0.18em}
    \item the project files;
    \item the commit history;
    \item the available branches;
    \item the remote connection to your GitHub fork.
\end{itemize}
\end{exampleblock}

\begin{alertblock}{Important}
Enter the command as one continuous terminal command. The line break
shown on the slide is only for visual readability.
\end{alertblock}

\end{frame}

% -----------------------------------------------------------------------------
\begin{frame}[fragile]{Alternative: Clone the Main Repository}
\framesubtitle{Use this only when instructed to obtain a read-only working copy}

\footnotesize

\begin{block}{Execute}
\begin{verbatim}
cd ~

git clone https://github.com/abdtnourji/\
quard_um6p_intership.git
\end{verbatim}
\end{block}

\begin{exampleblock}{When Is This Appropriate?}
This command is suitable when the objective is to obtain and run the
official project without pushing student modifications directly to it.
\end{exampleblock}

\begin{alertblock}{Contribution Limitation}
If you do not have write permission, you will not be able to push changes
to the supervisor's repository. For student contributions, use a fork or
the collaboration procedure specified by the instructor.
\end{alertblock}

\end{frame}

% -----------------------------------------------------------------------------
\begin{frame}[fragile]{Step 4: Enter the Repository}
\framesubtitle{Git commands must be executed from the correct project context}

\footnotesize

\begin{block}{Execute}
\begin{verbatim}
cd ~/quard_um6p_intership

pwd
ls -la
\end{verbatim}
\end{block}

\begin{exampleblock}{Expected Evidence}
The output of \texttt{pwd} should end with:

\begin{verbatim}
/quard_um6p_intership
\end{verbatim}

The output of \texttt{ls -la} should include the hidden
\texttt{.git} directory.
\end{exampleblock}

\begin{alertblock}{Common Mistake}
Cloning a repository does not automatically move the terminal into it.
You must enter the repository with \texttt{cd}.
\end{alertblock}

\end{frame}

% -----------------------------------------------------------------------------
\begin{frame}[fragile]{Step 5: Verify the Repository}
\framesubtitle{Confirm the branch, status, remote, and recent history}

\footnotesize

\begin{block}{Execute}
\begin{verbatim}
git status
git branch
git remote -v
git log --oneline -5
\end{verbatim}
\end{block}

\begin{exampleblock}{Expected Interpretation}
\begin{itemize}
    \setlength{\itemsep}{0.16em}
    \item Git recognizes the directory as a repository;
    \item the current branch is identified;
    \item \texttt{origin} points to the expected GitHub repository;
    \item recent commits are visible;
    \item the working tree is initially clean.
\end{itemize}
\end{exampleblock}

\begin{alertblock}{Stop If the Remote Is Incorrect}
Do not push until \texttt{git remote -v} shows the expected repository
owner and project name.
\end{alertblock}

\end{frame}

% -----------------------------------------------------------------------------
\begin{frame}[fragile]{Step 6: Read the Project Documentation}
\framesubtitle{Understand the repository before changing it}

\footnotesize

\begin{block}{Execute}
\begin{verbatim}
find . -maxdepth 2 -type d | sort
find . -name "README*"
less README.md
\end{verbatim}
\end{block}

\begin{exampleblock}{Inside the Installation Guide}
Locate and follow the installation files in this exact order:

\begin{enumerate}
    \setlength{\itemsep}{0.15em}
    \item \texttt{README\_PHASE0\_UBUNTU\_ROS2\_SETUP.md}
    \item \texttt{README\_PHASE\_1\_PX4\_AI\_VENV\_SETUP.md}
    \item \texttt{README\_PHASE\_2\_PIPELINE.md}
\end{enumerate}
\end{exampleblock}

\begin{alertblock}{Navigation Tip}
Press \texttt{q} to exit \texttt{less}. Do not begin Phase 1 before the
required Phase 0 verification checks have succeeded.
\end{alertblock}

\end{frame}

% -----------------------------------------------------------------------------
\begin{frame}[fragile]{Step 7: Synchronize Before Working}
\framesubtitle{Begin from the latest approved repository state}

\footnotesize

\begin{block}{First Identify the Current Branch}
\begin{verbatim}
git status
git branch
\end{verbatim}
\end{block}

\begin{block}{Then Pull the Current Main Branch}
\begin{verbatim}
git pull origin main
\end{verbatim}
\end{block}

\begin{exampleblock}{Possible Successful Output}
Git may report that the branch is already up to date, or it may download
and integrate newer commits.
\end{exampleblock}

\begin{alertblock}{Before Pulling}
The working tree should be clean and understood. If local modifications
already exist, inspect them before pulling.
\end{alertblock}

\end{frame}

% -----------------------------------------------------------------------------
\begin{frame}[fragile]{Step 8: Create a Practice Branch}
\framesubtitle{Keep the demonstration separate from the main branch}

\footnotesize

\begin{block}{Execute}
\begin{verbatim}
git switch -c docs/git-practice
\end{verbatim}
\end{block}

\begin{exampleblock}{Verify the Result}
\begin{verbatim}
git branch
git status
\end{verbatim}

The star displayed by \texttt{git branch} should appear next to:

\begin{verbatim}
docs/git-practice
\end{verbatim}
\end{exampleblock}

\begin{alertblock}{Why Create a Branch?}
The branch isolates the practice contribution from the stable project
history. It does not create another physical copy of the repository.
\end{alertblock}

\end{frame}

% -----------------------------------------------------------------------------
\begin{frame}[fragile]{Step 9: Create a Safe Student File}
\framesubtitle{Make a small documentation change without touching drone logic}

\footnotesize

\begin{block}{Execute}
\begin{verbatim}
mkdir -p student_work

echo "Student Git practice" > \
student_work/git_practice.txt
\end{verbatim}
\end{block}

\begin{exampleblock}{Verify the File}
\begin{verbatim}
cat student_work/git_practice.txt
git status
\end{verbatim}

Git should identify the file as \textbf{untracked}.
\end{exampleblock}

\begin{alertblock}{Engineering Boundary}
During this first exercise, do not modify PX4, Gazebo models, ROS 2
packages, launch files, or shared environment configuration.
\end{alertblock}

\end{frame}

% -----------------------------------------------------------------------------
\begin{frame}[fragile]{Step 10: Add Useful Project Evidence}
\framesubtitle{Document the environment used for the exercise}

\footnotesize

\begin{block}{Execute}
\begin{verbatim}
echo "Repository:" >> student_work/git_practice.txt
pwd >> student_work/git_practice.txt

echo "Current branch:" >> \
student_work/git_practice.txt

git branch --show-current >> \
student_work/git_practice.txt
\end{verbatim}
\end{block}

\begin{exampleblock}{Inspect the Result}
\begin{verbatim}
cat student_work/git_practice.txt
\end{verbatim}
\end{exampleblock}

\begin{alertblock}{Redirection Reminder}
The operator \texttt{>} replaces file content. The operator
\texttt{>>} appends new content without removing the existing lines.
\end{alertblock}

\end{frame}

% -----------------------------------------------------------------------------
\begin{frame}[fragile]{Step 11: Ask Git What Changed}
\framesubtitle{\texttt{git status} is the repository dashboard}

\footnotesize

\begin{block}{Execute}
\begin{verbatim}
git status
git status --short
\end{verbatim}
\end{block}

\begin{exampleblock}{Expected State}
The new file should appear as \textbf{untracked}. In the short format,
it may be marked with:

\begin{verbatim}
??
\end{verbatim}
\end{exampleblock}

\begin{alertblock}{Important Distinction}
The file exists on the computer, but Git is not tracking it yet.
Creating or saving a file does not automatically include it in a commit.
\end{alertblock}

\end{frame}

% -----------------------------------------------------------------------------
\begin{frame}[fragile]{Step 12: Stage the Intended File}
\framesubtitle{Select exactly what should enter the next commit}

\footnotesize

\begin{block}{Execute}
\begin{verbatim}
git add student_work/git_practice.txt
\end{verbatim}
\end{block}

\begin{exampleblock}{Verify the Staging Area}
\begin{verbatim}
git status
git diff --staged
\end{verbatim}

The file should now appear under the changes prepared for the next commit.
\end{exampleblock}

\begin{alertblock}{Why Use an Explicit Path?}
Avoid using \texttt{git add .} during the first exercises. It may stage
unrelated files, generated logs, large outputs, or sensitive information.
\end{alertblock}

\end{frame}

% -----------------------------------------------------------------------------
\begin{frame}[fragile]{Step 13: Review Before Committing}
\framesubtitle{The staging area is the proposed content of the next checkpoint}

\footnotesize

\begin{block}{Execute}
\begin{verbatim}
git diff --staged
\end{verbatim}
\end{block}

\begin{exampleblock}{Student Verification Questions}
Before committing, answer:

\begin{itemize}
    \setlength{\itemsep}{0.18em}
    \item Is this the intended file?
    \item Does it contain only appropriate information?
    \item Are there passwords, tokens, or private paths?
    \item Does the content support the objective of the commit?
\end{itemize}
\end{exampleblock}

\begin{alertblock}{Security Check}
Never commit passwords, access tokens, private keys, secret environment
variables, or confidential information.
\end{alertblock}

\end{frame}

% -----------------------------------------------------------------------------
\begin{frame}[fragile]{Step 14: Create the Commit}
\framesubtitle{Record one coherent and understandable project change}

\footnotesize

\begin{block}{Execute}
\begin{verbatim}
git commit -m \
"docs: add student Git practice evidence"
\end{verbatim}
\end{block}

\begin{exampleblock}{Why Is This a Good Message?}
\begin{itemize}
    \setlength{\itemsep}{0.18em}
    \item \texttt{docs} identifies the type of contribution;
    \item the remainder explains what was added;
    \item the message describes an engineering result.
\end{itemize}
\end{exampleblock}

\begin{alertblock}{Avoid Weak Messages}
Messages such as \texttt{update}, \texttt{final}, \texttt{changes}, or
\texttt{fix stuff} do not explain the purpose of the commit.
\end{alertblock}

\end{frame}

% -----------------------------------------------------------------------------
\begin{frame}[fragile]{Step 15: Verify the Local Commit}
\framesubtitle{A commit should create a clean and understandable local state}

\footnotesize

\begin{block}{Execute}
\begin{verbatim}
git status
git log --oneline -5
\end{verbatim}
\end{block}

\begin{exampleblock}{Expected Evidence}
\begin{itemize}
    \setlength{\itemsep}{0.18em}
    \item the working tree should be clean;
    \item the new commit should appear at the top of the history;
    \item the meaningful commit message should be visible.
\end{itemize}
\end{exampleblock}

\begin{alertblock}{Important Distinction}
The commit currently exists in the local repository. It has not yet been
published to GitHub.
\end{alertblock}

\end{frame}

% -----------------------------------------------------------------------------
\begin{frame}[fragile]{Step 16: Push the Practice Branch}
\framesubtitle{Publish the local branch to the student's GitHub fork}

\footnotesize

\begin{block}{Execute}
\begin{verbatim}
git push -u origin docs/git-practice
\end{verbatim}
\end{block}

\begin{exampleblock}{Meaning of the Command}
\begin{itemize}
    \setlength{\itemsep}{0.18em}
    \item \texttt{push} publishes local commits;
    \item \texttt{origin} identifies the remote repository;
    \item \texttt{docs/git-practice} identifies the branch;
    \item \texttt{-u} records the tracking relationship.
\end{itemize}
\end{exampleblock}

\begin{alertblock}{Permission Error}
If the push is rejected because you do not have permission, verify that
\texttt{origin} points to your fork and not directly to the supervisor's
repository.
\end{alertblock}

\end{frame}

% -----------------------------------------------------------------------------
\begin{frame}[fragile]{Step 17: Verify the Published Result}
\framesubtitle{Do not assume that a successful-looking command completed the workflow}

\footnotesize

\begin{block}{Verify in the Terminal}
\begin{verbatim}
git status
git branch
git remote -v
git log --oneline -3
\end{verbatim}
\end{block}

\begin{exampleblock}{Verify on GitHub}
Open your GitHub fork and confirm that:

\begin{itemize}
    \setlength{\itemsep}{0.18em}
    \item the branch \texttt{docs/git-practice} exists;
    \item the new commit is visible;
    \item the committed file is present;
    \item the commit message is correct.
\end{itemize}
\end{exampleblock}

\begin{alertblock}{Definition of Success}
The contribution is complete only when the local state and the GitHub
state agree and the published change can be inspected.
\end{alertblock}

\end{frame}

% -----------------------------------------------------------------------------
\begin{frame}{Optional Contribution Step}
\framesubtitle{Propose the branch without changing the main project directly}

\footnotesize

\begin{block}{Open a Pull Request}
On GitHub, select the pushed branch and create a pull request toward the
repository and branch specified by the instructor.
\end{block}

\begin{exampleblock}{A Good Pull Request Explains}
\begin{itemize}
    \setlength{\itemsep}{0.20em}
    \item what was changed;
    \item why the change was necessary;
    \item how the change was tested;
    \item which limitations or open questions remain.
\end{itemize}
\end{exampleblock}

\begin{alertblock}{Important}
Do not create unnecessary pull requests against the supervisor's
repository during a classroom exercise. Follow the contribution policy
defined for the internship.
\end{alertblock}

\end{frame}

% -----------------------------------------------------------------------------
\begin{frame}[fragile]{If the Wrong File Was Staged}
\framesubtitle{Remove it from staging without deleting the student's work}

\footnotesize

\begin{block}{Inspect the Current State}
\begin{verbatim}
git status
\end{verbatim}
\end{block}

\begin{exampleblock}{Remove One File from the Staging Area}
\begin{verbatim}
git restore --staged path/to/file
git status
\end{verbatim}
\end{exampleblock}

\begin{block}{What Happens?}
The file remains on the computer with its modifications, but it is no
longer selected for the next commit.
\end{block}

% \begin{alertblock}{Key Question}
% Do you want to remove the file from the next commit, or do you want to
% destroy its modification? These are different intentions.
% \end{alertblock}

\end{frame}

% -----------------------------------------------------------------------------
\begin{frame}[fragile]{If an Unwanted Edit Must Be Discarded}
\framesubtitle{Review the difference before removing uncommitted work}

\footnotesize

\begin{block}{First Inspect the Change}
\begin{verbatim}
git diff -- path/to/file
\end{verbatim}
\end{block}

\begin{exampleblock}{Discard the Unstaged Modification}
\begin{verbatim}
git restore path/to/file
git status
\end{verbatim}
\end{exampleblock}

\begin{alertblock}{Irreversible Local Change}
This operation removes the uncommitted modification from the selected
tracked file. Create a backup or use another branch if the experimental
work may still be valuable.
\end{alertblock}

\end{frame}

% -----------------------------------------------------------------------------
\begin{frame}[fragile]{If a Push Is Rejected}
\framesubtitle{Synchronize and understand the remote work instead of forcing}

\footnotesize

\begin{block}{First Collect Evidence}
\begin{verbatim}
git status
git branch
git remote -v
\end{verbatim}
\end{block}

\begin{exampleblock}{Then Synchronize the Intended Branch}
\begin{verbatim}
git pull origin main
\end{verbatim}

Read the complete output and resolve any conflict carefully before
attempting another push.
\end{exampleblock}

\begin{alertblock}{Never Use Force Automatically}
Do not respond to a rejected push with \texttt{git push --force}.
A forced push may overwrite shared history and remove another
contributor's work.
\end{alertblock}

\end{frame}

% -----------------------------------------------------------------------------
\begin{frame}[fragile]{Daily Workflow After the First Clone}
\framesubtitle{Do not clone the repository again every day}

\footnotesize

\begin{block}{Start a New Working Session}
\begin{verbatim}
cd ~/quard_um6p_intership
git status
git branch
git pull origin main
\end{verbatim}
\end{block}

\begin{exampleblock}{After Making a Focused Change}
\begin{verbatim}
git status
git diff
git add path/to/intended_file
git diff --staged
git commit -m "type: explain the change"
git push
\end{verbatim}
\end{exampleblock}

\begin{alertblock}{Mental Model}
Clone once. Pull to receive new work. Commit to preserve local progress.
Push to publish the commits to the configured remote branch.
\end{alertblock}

\end{frame}

% -----------------------------------------------------------------------------
\begin{frame}{Live Demonstration Summary}
\framesubtitle{Explain the state transition represented by each command}

\footnotesize

\begin{center}
\scriptsize
\texttt{git clone}
$\longrightarrow$
local repository

\vspace{0.12cm}

edit file
$\longrightarrow$
working-tree modification

\vspace{0.12cm}

\texttt{git add}
$\longrightarrow$
staged content

\vspace{0.12cm}

\texttt{git commit}
$\longrightarrow$
local history

\vspace{0.12cm}

\texttt{git push}
$\longrightarrow$
published GitHub branch
\end{center}

\pause

\begin{exampleblock}{Note}
A file is not published simply because it was saved. Git requires us to
review the change, select it, record it locally, and deliberately
synchronize it with GitHub.
\end{exampleblock}

\end{frame}

% -----------------------------------------------------------------------------
\begin{frame}{Student Challenge}
\framesubtitle{Repeat the complete workflow independently}

\footnotesize

\begin{block}{Mission}
Create a new branch named:

\begin{center}
\texttt{docs/installation-observations}
\end{center}

Then create a short file describing:

\begin{itemize}
    \setlength{\itemsep}{0.18em}
    \item the installed Ubuntu version;
    \item the current ROS 2 distribution;
    \item the installation phase currently completed;
    \item one verified command and its observed output.
\end{itemize}
\end{block}

\begin{exampleblock}{Expected Git Evidence}
The branch, staged diff, meaningful commit, and published GitHub branch
must all be visible and explainable.
\end{exampleblock}

% \begin{alertblock}{Security Requirement}
% Do not include passwords, access tokens, private keys, recovery keys,
% personal data, or complete secret environment files.
% \end{alertblock}

\end{frame}